\documentclass[11pt]{article}
\usepackage[a4paper,margin=2.6cm]{geometry}
\usepackage{amsmath,amssymb,amsthm}
\newtheorem{lemma}{Lemma}
\title{Landau's bounds on the early zeros\\
\large necessary density, local deficit, no constant $A$}
\author{D.\ Joubert}
\date{4 September 2026}
\begin{document}
\maketitle
\begin{abstract}
\noindent
Landau's necessary conditions say that a sampling set for
$PW_{L/2}$ must have lower density $D^-\ge L/(2\pi)$. Globally the
zeros satisfy this (the density tends to infinity). Locally, on
the desert and on $[0,2\pi\mu]$, they do not: the Landau dimension
of functions concentrated there exceeds the number of zeros by
$4$ to $90$, depending on the window. That deficit is why $c_L$
is exponentially small. It does not produce the constant $A$ in
the sampling inequality --- Landau is a necessary count of modes,
not a Slepian eigenvalue. Beurling's sufficiency is global and
too coarse to see the desert.
\end{abstract}

\section{What Landau proves}
For the Paley--Wiener space $PW_\Omega$ of type $\Omega$:
if $\Lambda$ is sampling then $D^-(\Lambda)\ge\Omega/\pi$;
if $\Lambda$ is interpolating then $D^+(\Lambda)\le\Omega/\pi$.
Here $\Omega=L/2$, so the critical density is
\[
\rho=\frac{L}{2\pi}=\frac1{2\pi/L},
\]
the reciprocal of the Nyquist spacing. The proof is spectral:
the concentration operator $\chi_I P_\Omega\chi_I$ has
\[
\#\{\text{eigenvalues close to }1\}
=\frac\Omega\pi\,|I|+o(|I|)
=\frac{L}{2\pi}\,|I|+o(|I|)
\]
as $|I|\to\infty$ (Landau 1967). Those eigenfunctions live on $I$
and are invisible to $\Lambda\setminus I$. If
$\#(\Lambda\cap I)$ is smaller than that dimension, $\Lambda$
cannot sample $PW_\Omega$.

\section{Global versus local}
The Riemann zeros (and those of a fixed primitive $L$-function)
have $N(T)\sim (T/2\pi)\log(T/2\pi)$. On $\mathbb{R}$ one has
$D^-=\infty>\rho$ for every finite $L$. Beurling's sufficiency,
applied to a separated subset as in \emph{sampling-floor}
Theorem~1, therefore gives $c_L>0$ under RH. The bound does not
see any desert: $D^-$ is a $\limsup$ of counts in long intervals,
and long intervals are well filled.

The desert is a \emph{short} interval. Landau's count on
$I=[-\gamma_1,\gamma_1]$ is finite and sharp enough to be used:
\[
\dim\{\text{functions concentrated on }[-\gamma_1,\gamma_1]\}
\approx\frac{L\gamma_1}{\pi}.
\]
One even pair $\{\pm\gamma_1\}$ is two points. The local density
$1/\gamma_1$ against $\rho$ is $0.16$--$0.40$ for $\zeta$ on the
windows below.

\section{Measured deficit}
Number of positive zeros up to $T$ versus Landau dimension
$LT/\pi$ (even window, $I=[-T,T]$).

\begin{center}
\small
\begin{tabular}{lrrrrrr}
set & $\mu$ & $T$ & need $LT/\pi$ & have & deficit & $D_{\mathrm{loc}}/\rho$\\
$\zeta$ & 3 & $\gamma_1$ & 4.9 & 1 & 3.9 & 0.40\\
$\zeta$ & 11 & $\gamma_1$ & 10.8 & 1 & 9.8 & 0.19\\
$\zeta$ & 11 & $2\pi\mu$ & 52.8 & 16 & 36.8 & 0.61\\
$\zeta$ & 16 & $2\pi\mu$ & 88.7 & 29 & 59.7 & 0.65\\
$\chi_5$ & 30 & $\gamma_1$ & 7.2 & 1 & 6.2 & 0.28\\
$\chi_5$ & 30 & $150$ & 161 & 89 & 72 & 1.11\\
$\chi_{17}$ & 11 & $\gamma_1$ & 2.9 & 1 & 1.9 & 0.70\\
\end{tabular}
\end{center}

On the desert the set is always subcritical (ratio $0.2$--$0.7$).
On $[0,2\pi\mu]$ it climbs toward criticality and crosses it for
$\chi_5$ once the cache is long enough --- that crossing is why a
finite list of zeros can look ``dense enough'' while $c_L$ is
still $e^{-s\mu}$: the \emph{first} $LT/\pi$ modes, not the last,
set the floor.

\section{What the bound does not give}
Landau does not estimate the gap $1-\lambda$ of the concentration
operator, only how many eigenvalues sit near $1$. The size of
$c_L$ is that gap for the ground state of the Gram of $\Lambda$,
which is a Slepian problem on a multi-band set (the complement of
the early gaps). Replacing the count ``deficit $=4$'' by
``$-\ln c_L=16.7$'' is exactly the passage this note does not
make --- it is \emph{desert-slepian}, and the factor is $2.5$--$6$.

Beurling's $D^->\rho$ is sufficient after thinning, and ignores
the same short interval. A quantitative Landau inequality with an
explicit $A$ in terms of $D^--\rho$ exists only in asymptotic
form for long sets; it does not cover a hole of length $\gamma_1$
followed by twenty sparse points.

\section*{Status}
Landau names the obstruction: on every window the early zeros are
a locally subcritical set, by $4$ to $90$ modes. That is the
necessary half of the sampling theorem already used qualitatively
in Theorem~1. It is not a formula for $c_L$. The sufficient
half with a constant remains the Slepian of the early set as a
whole. Not RH.
\end{document}
